%% ==================================================
%% Chapter: Preparation of matrix product states by log-depth circuits
%% Source: arXiv:2307.01696 (Malz, Styliaris, Wei, Cirac), main text and
%% Supplemental Material; the approximation estimate rests on
%% arXiv:2103.13367, Supplemental Material, eq. (S29).
%% Locate passages by equation label (eq:B_TM, eq:key_approximation,
%% eq:normal_fp_local, eq:phi_tilde, eq:unitary, eq:sequential_svds,
%% eq:C_is_an_isometry, eq:RG_circuit, eq:fp_general_ti,
%% eq:app_tidle_phi_1, eq:app_error, lm:1, lm:2, th:1).
%% ==================================================

\chapter{Preparation of Matrix Product States by Log-Depth Circuits}
\label{ch:log_depth_preparation}

A translation-invariant matrix product state on $N$ sites generated by a
normal tensor is approximated by blocking $q$ sites, replacing the positive
part of the polar decomposition of the blocked tensor by its renormalization
fixed point, and keeping the isometric part. The fixed point is a product of
nearest-neighbour entangled pairs, and the isometric part factorizes into
isometries with bounded bond dimension, either sequentially or along a binary
tree. With $q$ proportional to $\log(N/\varepsilon)$ this gives a circuit of
depth $O(\log(N/\varepsilon))$ with error $\varepsilon$, and no circuit of
depth $o(\log N)$ achieves a small error. The presentation follows
\cite{Malz2024Preparation}.

Throughout, $A$ is a tensor with physical dimension $d$ and bond dimension
$D$, $\E_A$ is its transfer map (Definition~\ref{def:transfer_map}), and
\begin{align}
    \ket{\phi_N(A)}
     & =\sum_{i_1,\dots,i_N=0}^{d-1}\Tr\left(A^{i_1}\cdots A^{i_N}\right)
    \ket{i_1\cdots i_N}
    \label{eq:ldp_ti_mps}
\end{align}
is the translation-invariant vector it generates, that is,
$\ket{\phi_N(A)}=\ket{V^{(N)}(A)}$ in the notation of
Definition~\ref{def:mpv_state}; the chapter keeps the source's notation
$\phi_N$. For two unit vectors the error is
$\varepsilon(\phi,\psi)=1-\lvert\braket{\phi}{\psi}\rvert$.

%% Convention: the source writes the transfer matrix as
%% E_A = sum_i conj(A^i) (x) A^i with right eigenvector |rho> and left
%% eigenvector <1|. In the transfer-map notation used here this reads
%% E_A(rho) = rho and sum_i (A^i)^dagger A^i = 1. The pair state below is
%% written in coordinates; the source display (1 (x) sqrt(rho)) sum_i |ii>
%% agrees with it after the transposition rho -> rho^T that the change of
%% vectorization convention induces.

\section{Blocking and polar decomposition}

\begin{definition}[Blocked map]
    \label{def:ldp_blocked_map}
    \uses{def:blocked_tensor}
    For $q\ge1$ the \emph{blocked map} of $A$ is the linear map
    $\mathcal B_q\colon\C^D\otimes\C^D\to(\C^d)^{\otimes q}$,
    \begin{align}
        \mathcal B_q\ket{\alpha,\beta}
         & =\sum_{i_1,\dots,i_q=0}^{d-1}
        \left(A^{i_1}\cdots A^{i_q}\right)_{\alpha\beta}\ket{i_1\cdots i_q},
        \label{eq:ldp_blocked_map}
    \end{align}
    that is, the $q$-blocked tensor read as a map from the virtual to the
    physical space \cite[eq.~(6) and the paragraph ``Approximation through the
        fixed-point state'']{Malz2024Preparation}.
\end{definition}

\begin{definition}[Tensor of a map on the doubled bond space]
    \label{def:ldp_tensor_of_map}
    A linear map $P$ on $\C^D\otimes\C^D$ defines a tensor with physical space
    $\C^D\otimes\C^D$ and bond dimension $D$ whose matrices are
    $(P^{(\gamma,\delta)})_{\alpha\beta}=\bra{\gamma,\delta}P\ket{\alpha,\beta}$.
    For a linear map $V\colon\C^D\otimes\C^D\to\C^{d'}$ the tensor $VP$ has
    matrices $(VP)^{j}=\sum_{\gamma,\delta}V_{j,(\gamma,\delta)}P^{(\gamma,\delta)}$.
\end{definition}

\begin{lemma}[Physical maps factor through the generated vector]
    \label{lem:ldp_physical_map_factorization}
    \uses{def:ldp_tensor_of_map, def:mpv_state}
    For every linear map $V\colon\C^D\otimes\C^D\to\C^{d'}$, every tensor $P$
    with physical space $\C^D\otimes\C^D$, and every $M\ge1$,
    \begin{align}
        \ket{\phi_M(VP)}
         & =V^{\otimes M}\ket{\phi_M(P)}.
        \label{eq:ldp_physical_map_factorization}
    \end{align}
\end{lemma}
\begin{proof}
    Each matrix of $VP$ is a linear combination of matrices of $P$ with
    coefficients given by $V$; expand the trace of the product by
    multilinearity.
\end{proof}

\begin{lemma}[Polar decomposition of the blocked map]
    \label{lem:ldp_polar_blocked_map}
    \uses{def:ldp_blocked_map}
    Let $P=(\mathcal B_q^\dagger\mathcal B_q)^{1/2}$ and let $\Pi$ be the
    orthogonal projection onto the range of $P$. There is a linear map
    $V\colon\C^D\otimes\C^D\to(\C^d)^{\otimes q}$ with
    \begin{align}
        \mathcal B_q
         & =VP,
         &
        V^\dagger V
         & =\Pi.
        \label{eq:ldp_polar_blocked_map}
    \end{align}
    If $\mathcal B_q$ is injective, then $P$ is positive definite and
    $V^\dagger V=\Id$. This is the polar decomposition used in
    \cite[eq.~(8) and Supplemental Material]{Malz2024Preparation}.
\end{lemma}

\begin{theorem}[The blocked state is an isometry applied to the positive part]
    \label{thm:ldp_blocked_state_factorization}
    \uses{lem:ldp_polar_blocked_map, lem:ldp_physical_map_factorization}
    With $V$ and $P$ as in Lemma~\ref{lem:ldp_polar_blocked_map}, for every
    $M\ge1$, grouping the $qM$ sites into $M$ consecutive blocks of $q$ sites,
    \begin{align}
        \ket{\phi_{qM}(A)}
         & =V^{\otimes M}\ket{\phi_M(P)},
         &
        \E_{A}^{\,q}
         & =\E_P.
        \label{eq:ldp_blocked_state_factorization}
    \end{align}
    This is the exact decomposition behind
    \cite[eq.~(8)]{Malz2024Preparation} and the display
    $\ket{\phi_N}=c_N^{-1}(\bigotimes_iV_i)\sum_j\beta_j\ket{v_{\mathrm{pos},j}}$
    of its Supplemental Material.
\end{theorem}
\begin{proof}
    The matrices of the $q$-blocked tensor are those of the tensor $VP$, so the
    first identity is Lemma~\ref{lem:ldp_physical_map_factorization}. For the
    second, $\E_A^{\,q}$ is the transfer map of the blocked tensor, and
    substituting $VP$ produces the factor $V^\dagger V=\Pi$, which acts
    trivially on the range of $P$.
\end{proof}

\section{The fixed-point pair state}

In this section $A$ is normal (Definition~\ref{def:normal}) and in the gauge
\begin{align}
    \sum_{i=0}^{d-1}(A^i)^\dagger A^i
     & =\Id,
     &
    \E_A(\rho)
     & =\rho,
     &
    \rho
     & >0,
     &
    \Tr\rho
     & =1,
    \label{eq:ldp_normal_gauge}
\end{align}
so that $\E_A^{\,n}(X)\to\Tr(X)\,\rho$ as $n\to\infty$
\cite[eq.~(5)]{Malz2024Preparation}. The limit uses normality as follows.
Normality means that the words of some fixed length $L_0$ span $\MN{D}$, so
by Theorem~\ref{thm:wielandt_proposition_three} the map $\E_A$ is strongly
irreducible: its peripheral spectrum is $\{1\}$ and its fixed point $\rho$
is unique up to scale. Since $\E_A$ is trace preserving, $1$ is a
semisimple eigenvalue, and all other eigenvalues have modulus less than
$1$. A periodic tensor is not normal, because no blocking of it is
injective.

\begin{definition}[Fixed-point tensor and pair state]
    \label{def:ldp_fixed_point_pair}
    \uses{def:ldp_tensor_of_map}
    The \emph{fixed-point tensor} $P_\infty$ has physical space
    $\C^D\otimes\C^D$ and matrices
    $P_\infty^{(\gamma,\delta)}=\sqrt\rho\,\ketbra{\gamma}{\delta}$. The
    \emph{pair state} on a right leg $R$ and a left leg $L$, each $\C^D$, is
    \begin{align}
        \ket{\omega}_{RL}
         & =\sum_{\gamma,\delta=0}^{D-1}(\sqrt\rho)_{\delta\gamma}
        \ket{\delta}_R\ket{\gamma}_L.
        \label{eq:ldp_pair_state}
    \end{align}
    These are the fixed-point tensor $P_\infty$ and the nearest-neighbour pair of
    \cite[eqs.~(8) and~(12)]{Malz2024Preparation}.
\end{definition}

\begin{lemma}[Transfer map of the fixed-point tensor]
    \label{lem:ldp_fixed_point_transfer}
    \uses{def:ldp_fixed_point_pair, def:mps_transfer_idempotence}
    For every $X\in\MN{D}$, $\E_{P_\infty}(X)=\Tr(X)\,\rho$. In particular
    $\E_{P_\infty}\circ\E_{P_\infty}=\E_{P_\infty}$, so $P_\infty$ is a
    renormalization fixed point in the sense of
    Definition~\ref{def:mps_transfer_idempotence}.
\end{lemma}
\begin{proof}
    $\sum_{\gamma,\delta}\sqrt\rho\ketbra{\gamma}{\delta}X\ketbra{\delta}{\gamma}\sqrt\rho
        =\Tr(X)\sqrt\rho\,\Id\,\sqrt\rho$, and $\Tr\rho=1$ gives idempotence.
\end{proof}

\begin{lemma}[Convergence of the positive part]
    \label{lem:ldp_positive_part_converges}
    \uses{lem:ldp_polar_blocked_map, def:ldp_fixed_point_pair}
    Let $P_q=(\mathcal B_q^\dagger\mathcal B_q)^{1/2}$. Then $P_q\to P_\infty$ as
    $q\to\infty$, where $P_\infty$ is read as the map
    $\ket{\alpha,\beta}\mapsto\sum_\gamma(\sqrt\rho)_{\alpha\gamma}\ket{\gamma,\beta}$.
    This is the limit in \cite[eq.~(8)]{Malz2024Preparation}.
\end{lemma}
\begin{proof}
    The matrix of $\mathcal B_q^\dagger\mathcal B_q$ is a rearrangement of
    $\E_A^{\,q}$, which converges to $X\mapsto\Tr(X)\rho$; the rearrangement of
    the limit is $\rho^{\mathsf T}\otimes\Id$, and the square root is continuous
    on positive semidefinite matrices.
\end{proof}

\begin{theorem}[The fixed-point state is a product of pairs]
    \label{thm:ldp_fixed_point_state_pairs}
    \uses{def:ldp_fixed_point_pair}
    For $M\ge1$ write the physical space of the $k$-th site of $P_\infty$ as
    $L_k\otimes R_k$. Then, with indices of $L$ read modulo $M$,
    \begin{align}
        \ket{\phi_M(P_\infty)}
         & =\bigotimes_{k=1}^M\ket{\omega}_{R_kL_{k+1}},
         &
        \braket{\omega}{\omega}
         & =1.
        \label{eq:ldp_fixed_point_state_pairs}
    \end{align}
    This is \cite[eq.~(12)]{Malz2024Preparation}.
\end{theorem}
\begin{proof}
    $\Tr\bigl(P_\infty^{(\gamma_1,\delta_1)}\cdots P_\infty^{(\gamma_M,\delta_M)}\bigr)
        =\prod_{k=1}^M(\sqrt\rho)_{\delta_k\gamma_{k+1}}$, and
    $\braket{\omega}{\omega}=\Tr\rho$.
\end{proof}

\begin{definition}[Approximating state]
    \label{def:ldp_approximating_state}
    \uses{lem:ldp_polar_blocked_map, def:ldp_fixed_point_pair}
    With $V$ from Lemma~\ref{lem:ldp_polar_blocked_map} and $N=qM$, the
    \emph{approximating state} is
    $\ket{\widetilde\phi_N}=\ket{\phi_M(VP_\infty)}/\lVert\phi_M(VP_\infty)\rVert$.
    By Lemma~\ref{lem:ldp_physical_map_factorization} and
    Theorem~\ref{thm:ldp_fixed_point_state_pairs},
    $\ket{\phi_M(VP_\infty)}=V^{\otimes M}\bigotimes_k\ket{\omega}_{R_kL_{k+1}}$.
    This is \cite[eqs.~(9) and~(10)]{Malz2024Preparation}.
\end{definition}

\begin{lemma}[A unitary implementing the isometry]
    \label{lem:ldp_unitary_implementing_isometry}
    \uses{lem:ldp_polar_blocked_map, lem:isometry_unitary_completion}
    Let $d\ge2$, put $r=\lceil\log_dD\rceil$, and suppose $q\ge2r$ and
    $\mathcal B_q$ is injective. Write
    $(\C^d)^{\otimes q}=L\otimes C\otimes R$, where $L$ and $R$ are the first
    and the last $r$ sites of the block and $C$ is the remaining $q-2r$ sites,
    embed $\C^D$ into $L$ and into $R$ as the span of the first $D$
    computational basis vectors, and let $\ket{0}_C=\ket{0}^{\otimes(q-2r)}$.
    Then there is a unitary $U$ on $(\C^d)^{\otimes q}$ with
    $U\bigl(\ket{\gamma}_L\ket{0}_C\ket{\delta}_R\bigr)=V\ket{\gamma,\delta}$
    for all $\gamma,\delta$, and
    \begin{align}
        \ket{\widetilde\phi_N}
         & =\Bigl(\bigotimes_{k=1}^M U_k\Bigr)
        \bigotimes_{k=1}^M\bigl(\ket{\omega}_{R_kL_{k+1}}\ket{0}_{C_k}\bigr).
        \label{eq:ldp_unitary_preparation}
    \end{align}
    This is \cite[eqs.~(10) and~(11)]{Malz2024Preparation}.
\end{lemma}
%% The source draws the circuit with D = d, where L and R are
%% single sites. For general D the legs are padded to r = ceil(log_d D)
%% sites each; the pair state then acts on 2r neighbouring sites.

\section{Exact factorizations of the isometry}

\begin{theorem}[Sequential factorization]
    \label{thm:ldp_sequential_factorization}
    \uses{lem:ldp_polar_blocked_map, lem:obc_left_isometric_representation}
    Suppose $\mathcal B_q$ is injective and let $V$ be as in
    Lemma~\ref{lem:ldp_polar_blocked_map}. There are integers
    $D'_1=D^2$, $D'_{q+1}=1$, and $D'_k\le D^2$, and isometries
    $W_k\colon\C^{D'_k}\to\C^d\otimes\C^{D'_{k+1}}$ for $k=1,\dots,q$, such that
    \begin{align}
        V
         & =\bigl(\Id_{d}^{\otimes(q-1)}\otimes W_q\bigr)\cdots
        \bigl(\Id_d\otimes W_2\bigr)W_1,
        \label{eq:ldp_sequential_factorization}
    \end{align}
    where $\C^{D'_1}$ is identified with $\C^D\otimes\C^D$ and the $k$-th
    tensor factor $\C^d$ is the $k$-th site of the block. This is
    \cite[eqs.~(13)--(15)]{Malz2024Preparation}.
\end{theorem}
\begin{proof}
    $V=\mathcal B_qP^{-1}$ is an open-boundary matrix product map with bond
    dimension at most $D^2$: the site matrices are $A^i\otimes\Id_D$, which carry
    the right virtual index along the block, and $P^{-1}$ is absorbed into the
    first site. Make the site maps $W_q,\dots,W_2$ isometric by the successive
    decompositions of Lemma~\ref{lem:obc_left_isometric_representation}. Then
    $V^\dagger V=W_1^\dagger W_1$, so $W_1$ is an isometry because $V$ is
    \cite[eq.~(15)]{Malz2024Preparation}.
\end{proof}

\begin{theorem}[Tree factorization]
    \label{thm:ldp_tree_factorization}
    \uses{lem:ldp_polar_blocked_map, lem:ldp_physical_map_factorization}
    Put $T_0=A$ and, for $j\ge1$, let $\mathcal B^{(j)}=V^{(j)}P_j$ be the polar
    decomposition of the two-site blocked map of $T_{j-1}$, and let $T_j$ be the
    tensor of $P_j$ (Definition~\ref{def:ldp_tensor_of_map}). Then for
    $q=2^k$,
    \begin{align}
        \mathcal B_q
         & =\Bigl(V^{(1)}\Bigr)^{\otimes2^{k-1}}
        \Bigl(V^{(2)}\Bigr)^{\otimes2^{k-2}}\cdots V^{(k)}\,P_k.
        \label{eq:ldp_tree_factorization}
    \end{align}
    Here $V^{(1)}$ maps $\C^{D^2}$ to $(\C^d)^{\otimes2}$ and each $V^{(j)}$ with
    $j\ge2$ maps $\C^{D^2}$ to $\C^{D^2}\otimes\C^{D^2}$. The product of the
    layers is a partial isometry with initial space the range of $P_k$, so by
    uniqueness of the polar decomposition it equals $V$ and $P_k=P_q$. This is
    \cite[eq.~(16)]{Malz2024Preparation}.
\end{theorem}
\begin{proof}
    Induction on $k$ using Lemma~\ref{lem:ldp_physical_map_factorization} for
    open products. The range of the two-site blocked map of $T_{j-1}$ lies in
    $\operatorname{ran}P_{j-1}\otimes\operatorname{ran}P_{j-1}$, so each layer
    is applied inside the initial space of the next.
\end{proof}

\section{Approximation error}

The correlation length of a normal tensor is $\xi=-1/\ln\lvert\lambda_2\rvert$,
where $\lambda_2$ is an eigenvalue of $\E_A$ of largest modulus other than
the Perron eigenvalue $1$. By the remark after (\ref{eq:ldp_normal_gauge}),
$\lvert\lambda_2\rvert<1$, so $\xi$ is finite.

\begin{lemma}[Overlap of the positive part with the fixed point]
    \label{lem:ldp_positive_part_overlap}
    \uses{def:ldp_fixed_point_pair, thm:ldp_blocked_state_factorization}
    Let $A$ be normal with correlation length $\xi$ and let $0<\gamma<1/2$.
    Then, for $N=qM$,
    \begin{align}
        \Bigl\lvert1-\bigl\lvert\bra{\phi_M(P_\infty)}\phi_M(P_q)\rangle\bigr\rvert\Bigr\rvert
         & =O\!\left(\frac Nq\,e^{-\gamma q/\xi}\right).
        \label{eq:ldp_positive_part_overlap}
    \end{align}
    This is \cite[Supplemental Material, eq.~(S29)]{Piroli2021LOCC}, quoted as
    \cite[eq.~(S9)]{Malz2024Preparation}.
\end{lemma}

\begin{theorem}[Approximation error, normal case]
    \label{thm:ldp_approximation_error_normal}
    \uses{lem:ldp_positive_part_overlap, def:ldp_approximating_state}
    Let $A$ be normal with correlation length $\xi$, let $\ket{\phi_N}$ be the
    normalized vector (\ref{eq:ldp_ti_mps}), and let $0<\gamma<1/2$. For
    $N=qM$ with $q\to\infty$ and $q=o(N)$,
    \begin{align}
        \varepsilon(\widetilde\phi_N,\phi_N)
         & =O\!\left(\frac Nq\,e^{-\gamma q/\xi}\right).
        \label{eq:ldp_approximation_error_normal}
    \end{align}
    In particular $q=\lceil\xi(1+\eta)\ln N/\gamma\rceil$ gives
    $e^{-\gamma q/\xi}\le N^{-(1+\eta)}$ and hence
    $\varepsilon=O(N^{-\eta})$ for every $\eta>0$. This is
    \cite[Lemma~1 and Lemma~1$'$(i)]{Malz2024Preparation}.
\end{theorem}
%% The source states Lemma 1'(i) without q = o(N); its proof uses q = o(N) to
%% absorb the O(exp(-N/xi)) normalization term. The statement above carries
%% the hypothesis the proof uses. The block length is scaled by 1/gamma
%% because the estimate is only available for gamma < 1/2, so the endpoint
%% choice q = 2 xi (1 + eta) ln N does not follow from it.

\section{Depth of the preparation}

\begin{definition}[Local circuit of depth $T$]
    \label{def:ldp_local_circuit}
    A \emph{local circuit of depth $T$} on a chain of $N$ sites is a product of
    $T$ layers, each a tensor product of unitaries acting on disjoint pairs of
    neighbouring sites. A state is \emph{prepared in depth $T$} if it equals a
    local circuit of depth $T$ applied to a product state.
\end{definition}

\begin{theorem}[Preparation in depth $O(\log(N/\varepsilon))$]
    \label{thm:ldp_depth_upper_bound}
    \uses{def:ldp_local_circuit, lem:ldp_unitary_implementing_isometry,
        thm:ldp_sequential_factorization, thm:ldp_tree_factorization,
        thm:ldp_approximation_error_normal}
    Let $A$ be normal. There is a constant $c$ depending only on $A$ such that,
    for every $N\ge2$ and $0<\varepsilon\le1$, a state $\ket{\psi_N}$ with
    $\varepsilon(\psi_N,\phi_N)\le\varepsilon$ is prepared in depth at most
    $c\log(N/\varepsilon)$. This is \cite[eq.~(1)]{Malz2024Preparation}.
\end{theorem}
%% The source illustrates the circuit with D = d. For general D the legs L_k,
%% R_k are realized on groups of ceil(log_d D) sites as in
%% Lemma lem:ldp_unitary_implementing_isometry; the depth estimate is
%% unchanged up to constants. The pair layer has constant depth, and the
%% sequential factorization with SWAP gates has depth O(q). The range
%% 0 < eps <= 1, N >= 2 keeps log(N/eps) positive.

\begin{lemma}[Exponentially decaying correlations]
    \label{lem:ldp_decaying_correlations}
    \uses{def:injective}
    Let $A$ be injective with correlation length $\xi>0$, and let
    $\ket{\phi_N}$ be the normalized vectors (\ref{eq:ldp_ti_mps}). There are
    operators $\mathcal O,\mathcal O'$ on $\C^d$ of norm one and $c>0$ such that
    for every $s>1$ and all sufficiently large $N$,
    \begin{align}
        \bra{\phi_N}\mathcal O_1\ket{\phi_N}
         & =\bra{\phi_N}\mathcal O'_s\ket{\phi_N}=0,
         &
        \bigl\lvert\bra{\phi_N}\mathcal O_1\mathcal O'_s\ket{\phi_N}\bigr\rvert
         & \ge c\,e^{-(s-1)/\xi}.
        \label{eq:ldp_decaying_correlations}
    \end{align}
    This is \cite[Lemma~2]{Malz2024Preparation}.
\end{lemma}
%% The bound is stated for the modulus: the connected correlation is
%% approximately lambda_2^(s-1) c', which is complex in general. To check
%% against the source proof: when several eigenvalues share the modulus
%% |lambda_2| their contributions can cancel at individual s, and the lower
%% bound then holds only along a subsequence of s, which suffices for
%% Theorem thm:ldp_depth_lower_bound.

\begin{theorem}[No preparation in depth $o(\log N)$]
    \label{thm:ldp_depth_lower_bound}
    \uses{def:ldp_local_circuit, lem:ldp_decaying_correlations,
        thm:ldp_approximation_error_normal}
    Let $A$ be normal with correlation length $\xi>0$, let $\ket{\phi_N}$ be the
    normalized vectors (\ref{eq:ldp_ti_mps}), and let $\ket{\psi_N}$ be prepared
    in depth $T_N$ with $T_N=o(\log N)$. Then there is $N_0$ such that
    $\varepsilon(\phi_N,\psi_N)>1/2$ for all $N>N_0$. This is
    \cite[Theorem~1]{Malz2024Preparation}.
\end{theorem}
%% The source proof applies Lemma 2, stated for injective A, to a normal A.
%% A normal tensor is injective after blocking a bounded number of sites; the
%% proof must be run on the blocked chain or Lemma 2 extended to normal A.

\section{Tensors that are not normal}

\begin{definition}[Fixed point of a basis of normal tensors]
    \label{def:ldp_non_normal_fixed_point}
    \uses{def:bnt, def:ldp_fixed_point_pair, lem:ldp_polar_blocked_map}
    Let $A^i=\bigoplus_{j=1}^b\operatorname{diag}(\mu_{j,1},\dots,\mu_{j,m_j})
        \otimes A_j^i$ with normal $A_j$, each in the gauge
    (\ref{eq:ldp_normal_gauge}), forming a basis of normal tensors, and put
    $\beta_j=\sum_{k=1}^{m_j}\mu_{j,k}^N$. Consider only lengths $N=qM$ with
    $(\beta_1,\dots,\beta_b)\neq0$; for the other lengths
    $\ket{\phi_N(A)}=\sum_j\beta_j\ket{\phi_N(A_j)}=0$ and there is nothing to
    prepare. With $\ket{\omega_j}$ the pair state of $A_j$ and
    $\alpha^{(N)}_j=\beta_j/(\sum_l\lvert\beta_l\rvert^2)^{1/2}$, the
    \emph{fixed-point state} is
    \begin{align}
        \ket{\Omega'}
         & =\sum_{j=1}^b\alpha^{(N)}_j\bigotimes_{k=1}^{M}\ket{\omega_j}_{R_kL_{k+1}}.
        \label{eq:ldp_non_normal_fixed_point}
    \end{align}
    With $V$ from Lemma~\ref{lem:ldp_polar_blocked_map} for the blocked map of
    $A$, the \emph{approximating state} is
    $\ket{\widetilde\phi_N}=V^{\otimes M}\ket{\Omega'}/
        \lVert V^{\otimes M}\ket{\Omega'}\rVert$, whenever the denominator is
    nonzero.
    This is \cite[eqs.~(19) and~(S7)]{Malz2024Preparation}, with the
    canonical form of \cite{Cirac2017MPDO_AnnPhys}.
\end{definition}

\begin{lemma}[GHZ form of the fixed-point state]
    \label{lem:ldp_ghz_form}
    \uses{def:ldp_non_normal_fixed_point}
    If $\braket{\omega_j}{\omega_{j'}}=\delta_{jj'}$, then
    $W\colon\ket{j}\mapsto\ket{\omega_j}$ is an isometry and
    $\ket{\Omega'}=W^{\otimes M}\ket{\chi_M}$ with
    $\ket{\chi_M}=\sum_j\alpha^{(N)}_j\ket{j}^{\otimes M}$, up to the
    regrouping of the legs $R_kL_{k+1}$.
    This is the preparation of \cite[eq.~(19)]{Malz2024Preparation}.
\end{lemma}

\begin{theorem}[Approximation error, general case]
    \label{thm:ldp_approximation_error_general}
    \uses{def:ldp_non_normal_fixed_point}
    With the notation of Definition~\ref{def:ldp_non_normal_fixed_point},
    including the approximating state $\widetilde\phi_N$ defined there,
    $\xi_{\mathrm{diag}}=\max_j\xi_{jj}$ the largest correlation length of the
    $A_j$, $0<\gamma<1/2$, and $q=o(N)$,
    \begin{align}
        \varepsilon(\widetilde\phi_N,\phi_N)
         & =O\!\left(\frac Nq\,e^{-\gamma q/\xi_{\mathrm{diag}}}\right).
        \label{eq:ldp_approximation_error_general}
    \end{align}
    This is \cite[Lemma~1$'$(ii)]{Malz2024Preparation}.
\end{theorem}
%% The direct sum A is not normal, so Lemma lem:ldp_positive_part_overlap
%% does not apply to it. The mixed-block overlap estimate is the content of
%% Lemma 1'(ii) of the source and is cited, not derived, here.
